\setcounter{section}{2} % Set to appropriate section number for Week 3

\section{Compactness (Continued)}

\begin{theorem*}[Heine-Borel]
A subset $K \subseteq \R^n$ equipped with the \emph{standard metric} is compact if and only if it is closed and bounded.
\end{theorem*}

\gemininote{The Heine-Borel theorem perfectly captures our intuitive visual understanding of a "compact" (small, contained, boundary-inclusive) set in standard geometry. However, this beautiful equivalence strictly relies on the standard Euclidean space. We will now prove that Heine-Borel completely breaks down if we use a non-standard metric!}

\begin{exercise*}
Let $(X, d)$ be a metric space.
\begin{enumerate}
    \item Show that $\tilde{d}(x,y) := \frac{d(x,y)}{1 + d(x,y)}$ is also a valid metric on $X$.
    \item Show that if $U \subseteq X$ is open with respect to $d$, then $U$ is also open with respect to $\tilde{d}$. \textcolor{eGray}{(Optional)}
    \item Show that for $X = \R^n$ with the metric $\tilde{d}(x,y) = \frac{\abs{x-y}}{1+\abs{x-y}}$, the entire space $\R^n \subseteq X$ is closed and bounded, but \emph{not} compact.
\end{enumerate}
\end{exercise*}

\begin{solution}
\textbf{a)} Non-negativity, definiteness, and symmetry follow immediately from the properties of $d(x,y)$. 
For the triangle inequality, we want to show:
\[ \frac{d(x,y)}{1+d(x,y)} \le \frac{d(x,z)}{1+d(x,z)} + \frac{d(z,y)}{1+d(z,y)} \]
Let's denote $d_1 = d(x,y)$, $d_2 = d(x,z)$, and $d_3 = d(z,y)$. We cross-multiply to clear the denominators:
\[ d_1(1+d_2)(1+d_3) \le d_2(1+d_1)(1+d_3) + d_3(1+d_1)(1+d_2) \]
Expanding everything yields:
\[ d_1 + d_1 d_2 + d_1 d_3 + d_1 d_2 d_3 \le d_2 + d_1 d_2 + d_2 d_3 + d_1 d_2 d_3 + d_3 + d_1 d_3 + d_2 d_3 + d_1 d_2 d_3 \]
Canceling the common terms ($d_1 d_2$, $d_1 d_3$, and one $d_1 d_2 d_3$) leaves us with:
\[ d_1 \le d_2 + 2d_2 d_3 + d_3 + d_1 d_2 d_3 \]
We know from the standard triangle inequality that $d_1 \le d_2 + d_3$. Subtracting this known truth from both sides yields:
\[ 0 \le 2d_2 d_3 + d_1 d_2 d_3 \]
Since distances are non-negative, this final inequality is trivially true! Thus, $\tilde{d}$ is a metric.

\vspace{0.3cm}
\textbf{b)} \textcolor{eGray}{(Proof sketch)} The function $t \mapsto \frac{t}{1+t}$ is strictly monotonically increasing for $t \ge 0$. Therefore, open balls in $d$ can be directly mapped to open balls in $\tilde{d}$ by simply transforming the radius $r$ into $\tilde{r} = \frac{r}{1+r}$.

\vspace{0.3cm}
\textbf{c)} Note that for any $t \ge 0$, the expression $\frac{t}{1+t}$ is strictly strictly less than $1$. Therefore, for all $x,y \in \R^n$:
\[ \tilde{d}(x,y) = \frac{\abs{x-y}}{1+\abs{x-y}} < 1 \]
This means the entire space $\R^n$ fits inside the open ball $\tilde{B}_1(0)$. Thus, $\R^n$ is \textbf{bounded} with respect to this metric! Also, the full space is trivially \textbf{closed}. 
However, consider the sequence $x_n = (n, 0, \dots, 0)$. This sequence wanders off to infinity and has absolutely no convergent subsequence. Since it lacks sequential compactness, $\R^n$ is \textbf{not compact}, breaking Heine-Borel!
\end{solution}

\hrulefill

\section{Why do we care about compactness?}

Compactness is one of the most powerful properties a space can have. It grants us three massive mathematical guarantees:

\begin{enumerate}
    \item \textbf{Preserved by continuous functions:} 
    Let $X, Y$ be metric spaces, and $f: X \to Y$ continuous. If $K \subseteq X$ is compact, then its image $f(K) \subseteq Y$ is guaranteed to be compact.
    
    \item \textbf{Weierstrass Theorem / Extreme Value Theorem:}
    Let $K \subseteq X$ be a compact metric space, and $f: K \to \R$ a continuous function. Then $f$ attains an absolute maximum and minimum within $K$. \textcolor{eGray}{(In particular, $f$ is bounded!)}
    
    \item \textbf{Uniform continuity:}
    Let $X, Y$ be metric spaces, and $f: X \to Y$ continuous. If $X$ is compact, then $f$ is automatically \emph{uniformly} continuous! \textcolor{eGray}{(You can pick a global $\delta$ that works everywhere for a given $\eps$.)}
\end{enumerate}

\begin{center}
\begin{tikzpicture}[scale=1.5, thick]
    % Axes
    \draw[->, eBlack] (-0.5, 0) -- (4, 0) node[right] {$x$};
    \draw[->, eBlack] (0, -0.5) -- (0, 3) node[above] {$f(x)$};
    
    % Function curve
    \draw[eMediumBlue, very thick] (0.5, 1.5) .. controls (1, 2.5) and (1.5, 3) .. (2.5, 0.5) .. controls (3, 0) and (3.5, 1) .. (3.8, 1.5);
    
    % Interval markers
    \draw[dashed, eGray] (0.5, 0) node[below, eBlack] {$a$} -- (0.5, 1.5);
    \draw[dashed, eGray] (3.8, 0) node[below, eBlack] {$b$} -- (3.8, 1.5);
    
    % Max / Min markers
    \coordinate (Max) at (1.15, 2.55);
    \coordinate (Min) at (2.7, 0.35);
    
    \fill[eDarkRed] (Max) circle (2pt);
    \fill[eDarkCyan] (Min) circle (2pt);
    
    \draw[dashed, eGray] (Max) -- (1.15, 0) node[below, eBlack] {$c$};
    \draw[dashed, eGray] (Max) -- (0, 2.55) node[left, eBlack] {$f(c)$};
    
    \draw[dashed, eGray] (Min) -- (2.7, 0) node[below, eBlack] {$d$};
    \draw[dashed, eGray] (Min) -- (0, 0.35) node[left, eBlack] {$f(d)$};
\end{tikzpicture}\\
\textcolor{eGray}{A continuous function on a compact interval $[a,b]$ always attains its absolute maximum (red) and absolute minimum (cyan).}
\end{center}

\hrulefill

\section{Connectedness}

\begin{definition*}[Disconnected metric space/subset]
A metric space $(X, d)$ is \textbf{disconnected} if there exist subsets $U_1, U_2 \subseteq X$ which are:
\begin{enumerate}
    \item \textbf{open}
    \item \textbf{non-empty} ($U_1 \ne \emptyset \ne U_2$)
    \item \textbf{disjoint} ($U_1 \cap U_2 = \emptyset$)
\end{enumerate}
such that $X = U_1 \cup U_2$. Similarly, a subset $A \subseteq X$ is disconnected if it can be split in this exact same way using the induced metric. \\
\textbf{Connected} is simply defined as the logical negation of disconnected!
\end{definition*}

\gemininote{Intuitively, a space is disconnected if it consists of two or more completely isolated \qt{islands}. Because proving connectedness directly is very hard, we almost always prove it by contradiction: We assume the space \emph{is} disconnected (split into two islands $U_1$ and $U_2$) and show that this leads to a logical paradox.}

\begin{center}
\begin{tikzpicture}[scale=1.2, very thick]
    % Disconnected
    \node[eDarkRed, font=\bfseries] at (1.5, 1.5) {Disconnected};
    \draw[eMediumBlue, fill=eMediumBlue!10] (0,0) circle (0.8cm);
    \draw[eMediumBlue, fill=eMediumBlue!10] (3,0) circle (1cm);
    \node[eMediumBlue] at (0,0) {$U_1$};
    \node[eMediumBlue] at (3,0) {$U_2$};
    \draw[<->, eDarkRed, dashed] (0.9,0) -- (1.9,0) node[midway, above] {Gap};

    % Connected
    \begin{scope}[xshift=6cm]
        \node[eDarkGreen, font=\bfseries] at (1.5, 1.5) {Connected};
        \draw[eMediumBlue, fill=eMediumBlue!10] (0,0) .. controls (1,1) and (2,1) .. (3,0) .. controls (2,-1) and (1,-1) .. (0,0);
        \node[eMediumBlue] at (1.5,0) {$X$};
    \end{scope}
\end{tikzpicture}
\end{center}

\begin{theorem*}
Continuous functions preserve connected sets! Let $X, Y$ be metric spaces, and $f: X \to Y$ continuous. If $A \subseteq X$ is connected, then its image $f(A)$ is strictly connected in $Y$.
\end{theorem*}

\begin{exercise*}
Are the following statements True or False?
\begin{enumerate}
    \item The union of connected subsets with a common point is connected.
    \item The intersection of connected subsets is connected.
    \item The closure of a connected subset is connected.
\end{enumerate}
\end{exercise*}

\begin{solution}
\textbf{1. True.} \\
\textcolor{eSaddleBrown}{Proof:} Let $X$ be a metric space, $V_1, V_2 \subseteq X$ connected, and let $x_0 \in V_1 \cap V_2$ be their shared point. Assume by contradiction that their union is disconnected. Then $V_1 \cup V_2 = U_1 \cup U_2$, where $U_1, U_2 \subseteq X$ are open, non-empty, and disjoint. 
Without loss of generality, assume the shared point $x_0$ lies in $U_1$. Because $x_0$ is in both $V_1$ and $V_2$, it follows that $U_1 \cap V_1 \ne \emptyset$ and $U_1 \cap V_2 \ne \emptyset$. 
Since $V_1$ and $V_2$ are individually connected, they cannot be split across $U_1$ and $U_2$. Since they both share a point with $U_1$, they must both lie \emph{entirely} inside $U_1$. Thus $V_1 \cap U_2 = \emptyset$ and $V_2 \cap U_2 = \emptyset$. This forces $U_2 = \emptyset$, which violates the definition of disconnectedness! Contradiction.

\vspace{0.3cm}
\textbf{2. False.} \\
\textcolor{eSaddleBrown}{Counterexample:} Let $X = \R^2$. Let $U = [-1, 1] \times \set{0}$ (a horizontal line segment) and $V = S^1$ (the unit circle). Both are perfectly connected. However, their intersection is exactly two isolated points: $U \cap V = \set{(-1,0), (1,0)}$, which is clearly disconnected!

\vspace{0.3cm}
\textbf{3. True.} \\
\textcolor{eSaddleBrown}{Proof:} Let $A$ be connected with closure $\bar{A}$. Suppose $\bar{A}$ is disconnected, meaning there exist disjoint open sets $U_1, U_2$ covering $\bar{A}$. Without loss of generality, let $A \cap U_1 \ne \emptyset$. Since $A$ is connected and covered by $U_1 \cup U_2$, it must lie entirely in one of them, meaning $A \cap U_2 = \emptyset$. 
However, $U_2$ is an open set. If an open set contains no points of $A$, it cannot possibly contain any boundary points of $A$ either! Thus $U_2 \cap \bar{A} = \emptyset$, meaning $U_2$ is empty. Contradiction.
\end{solution}

\hrulefill

\section{Path Connectedness}

\begin{definition*}[Paths and Path-connectedness]
Let $(X, d)$ be a metric space. 
\begin{itemize}
    \item A \textbf{path} is a continuous function $\gamma: [0,1] \to X$, where $[0,1]$ carries the standard metric.
    \item A space (or subset) $X$ is \textbf{path-connected} if for any two points $x, y \in X$, there exists a path $\gamma$ such that $\gamma(0) = x$ and $\gamma(1) = y$.
\end{itemize}
\end{definition*}

\begin{proposition*}
\begin{itemize}
    \item \textbf{Path-connected $\implies$ connected.} (This is always true!)
    \item For \emph{open} subsets of Euclidean space $\R^n$: \textbf{connected $\iff$ path-connected.}
\end{itemize}
\end{proposition*}

\subsection*{\warning{Counterexample: The Topologist's Sine Curve}}
The reverse implication (connected $\implies$ path-connected) is \textbf{false} for non-open sets! Consider the set:
\[ S := \set{\left(t, \sin\left(\frac{1}{t}\right)\right) : t > 0} \cup \set{(0,0)} \subseteq \R^2 \]
$S$ is \textbf{connected} (because the sine wave part is connected, and we proved earlier that adding boundary points via the closure maintains connectedness). However, $S$ is \textbf{not path-connected}!

\gemininote{Why can't we draw a path? As $t$ approaches $0$, the frequency of the sine wave goes to infinity. It oscillates up and down infinitely fast between $-1$ and $1$. If you try to draw a path from the origin $(0,0)$ to the rest of the curve, your pen would have to cross the x-axis an infinite number of times in a finite amount of time, shattering continuity!}

\begin{solution}
Assume by contradiction that $S$ is path-connected. Then there exists a continuous path $\gamma: [0,1] \to S$ starting at the origin $\gamma(0) = (0,0)$ and ending at some point on the curve, say $\gamma(1) = (1, \sin(1))$. 
We can write the path as two coordinate functions $\gamma(t) = (x(t), y(t))$. Since $\gamma$ is continuous, $x(t)$ and $y(t)$ must be continuous.

By the Intermediate Value Theorem, since $x(0) = 0$ and $x(1) = 1$, the function $x(t)$ must hit every x-coordinate in between. Let's look at the peaks of the sine wave, which occur at $x$-coordinates $\frac{1}{2\pi n + \pi/2}$. 
We can construct a sequence of times $(a_n)_{n=0}^\infty$ in $[0,1]$ such that:
\[ x(a_n) = \frac{1}{2\pi n + \pi/2} \]
Notice that the y-coordinate at all these peaks is exactly $y(a_n) = 1$. 
As $n \to \infty$, the x-coordinate goes to 0, so $a_n \to 0$. But look at the limits:
\[ \lim_{n \to \infty} \gamma(a_n) = (0, 1) \]
However, sequential continuity demands that $\lim \gamma(a_n) = \gamma(0) = (0,0)$. 
$(0,1) \ne (0,0)$, which is a massive contradiction! Thus, no such continuous path can exist.
\end{solution}

\hrulefill

\section{Inner Product and Norm}

\begin{definition*}[Inner Product]
Let $V$ be a real vector space. An \textbf{inner product} on $V$ is a map $\langle \cdot, \cdot \rangle: V \times V \to \R$ satisfying:
\begin{enumerate}
    \item \textbf{Bilinearity:} $\langle \alpha u + \beta v, w \rangle = \alpha \langle u, w \rangle + \beta \langle v, w \rangle$ \textcolor{eGray}{(and equivalently in the second variable)}
    \item \textbf{Symmetry:} $\langle u, v \rangle = \langle v, u \rangle$
    \item \textbf{Definiteness:} $\langle u, u \rangle \ge 0$, and $\langle u, u \rangle = 0 \iff u = 0$
\end{enumerate}
\end{definition*}

\begin{definition*}[Norm]
A \textbf{norm} on $V$ is a map $\norm{\cdot}: V \to [0, \infty)$ satisfying:
\begin{enumerate}
    \item \textbf{Definiteness:} $\norm{v} \ge 0$, and $\norm{v} = 0 \iff v = 0$
    \item \textbf{Homogeneity:} $\norm{\alpha v} = \abs{\alpha} \norm{v}$
    \item \textbf{Triangle Inequality:} $\norm{u + v} \le \norm{u} + \norm{v}$
\end{enumerate}
\end{definition*}

\subsection*{Geometric Meaning and Projections}
Every inner product naturally defines a norm via $\norm{v} := \sqrt{\langle v, v \rangle}$. 
In $\R^2$, the standard inner product geometrically yields the angle $\theta$ between two vectors:
\[ \langle x, y \rangle = \norm{x} \norm{y} \cos\theta \]
In a general inner product space, this formula actually serves as the \emph{definition} of an angle! 

This allows us to define the orthogonal projection of a vector $u$ onto a vector $v$:
\[ \pi_v(u) = \frac{\langle u, v \rangle}{\norm{v}^2} v \]
In words, $\pi_v(u)$ is the "shadow" of $u$ pointing exactly along the direction of $v$.

\begin{center}
\begin{tikzpicture}[scale=1.5, very thick]
    % Axes
    \draw[->, eBlack, thick] (0,0) -- (5,0);
    \draw[->, eBlack, thick] (0,0) -- (0,3);
    
    % Vectors
    \coordinate (V) at (4.5, 1.5);
    \coordinate (U) at (2.5, 2.8);
    \coordinate (Proj) at (3.3, 1.1); % Approx projection
    
    \draw[->, eSaddleBrown] (0,0) -- node[below right] {$v$} (V);
    \draw[->, eMediumBlue] (0,0) -- node[above left] {$u$} (U);
    
    % Projection elements
    \draw[->, eMediumOrchid] (0,0) -- node[below right, yshift=-0.1cm] {$\pi_v(u)$} (Proj);
    \draw[dashed, eMediumOrchid] (U) -- (Proj);
    
    % Angle
    \draw[eDarkGreen] (1,0.33) arc (18.4:48.2:1.05cm);
    \node[eDarkGreen] at (1.1, 0.7) {$\theta$};
    
    % Right angle square
    \draw[eBlack, thin] (3.15, 1.35) -- (3.35, 1.45) -- (3.45, 1.25);
    \fill[eBlack] (3.3, 1.35) circle (0.5pt);
\end{tikzpicture}
\end{center}

\subsection*{The Cauchy-Schwarz Inequality}
If we take the norm of the projection, we get:
\[ \norm{\pi_v(u)} = \norm{u} \abs{\cos\theta} = \frac{\abs{\langle u, v \rangle}}{\norm{v}} \]
Because the "shadow" of a vector can never be longer than the vector itself ($\abs{\cos\theta} \le 1$), it must hold that $\norm{\pi_v(u)} \le \norm{u}$. Moving the denominator over yields the legendary Cauchy-Schwarz Inequality:
\[ \abs{\langle u, v \rangle} \le \norm{u} \norm{v} \]

\gemininote{Hint for a rigorous algebraic proof: Use the definiteness property! We know that the length of the \qt{perpendicular part} of the vector must be $\ge 0$. So, expand the inner product of the perpendicular part with itself: $\langle u - \pi_v(u), u - \pi_v(u) \rangle \ge 0$. Expanding this term will elegantly spit out the Cauchy-Schwarz inequality!}
